\subsection{The full abstract length theorem}
\label{subsec:selmer-core-wrapper}

The preceding numerical comparison has four additional requirements:
a Smith basis for the negative equations, a chart without weight-zero
variables, a chart-nonzero representative of $a$, and a finite orbit
cover. Each follows from the original model hypotheses.

Put $I=\ker x$ and
$I^{(2)}=\{h:\exists\,0\ne c\in\mathcal O,\ ch\in I^2\}$, and
write $T=I/I^{(2)}$. Assume the original graded chart
\[
 \Phi:A\longrightarrow F[[z_1,\ldots,z_m,y_1,\ldots,y_n]]
\]
has constant coefficient $x$ and induces an isomorphism from
$T\otimes F$ to its space of linear forms, preserving weights.
The $z_i$ have weight zero; the numbers of $y$ variables of weights
$-1,-2,1,2$ are respectively $r,s,r'+1,s'+1$, with $r\ge1$.
Assume the classes of $F_i,G_j,Z_i$ span the rational cotangent
spaces of weights $-1,-2,0$, respectively. Retain the finite quotient,
fixed-point separation, two-sided tangent bound, and trace
nonvanishing hypotheses used above.

\begin{theorem}[Abstract length bound]
\label{thm:selmer-abstract-length}
The two-sided points of $A/Q$ have finitely many
$\mathbb G_m(K)$-orbits, in bijection with their restrictions to $B$.
The module
$\mathcal C^-=T_{-1}/\langle[F_i]\rangle_{\mathcal O}$ has finite
length. For every finite set $U$ of two-sided points meeting every
orbit,
\[
 \ell_{\mathcal O}\mathcal C^-
 \le2\sum_{\xi\in U}v_\pi(\xi(a)).
\]
\end{theorem}

\begin{proof}
Since $A$ is Noetherian, $I/I^2$ is finite over
$A/I=\mathcal O$. The ideal $I^{(2)}$ is the inverse image of its
$\mathcal O$-torsion, and lies in $I$. Thus $T$ is finite free,
as are its weight components. The chart identifies
$T_{-1}\otimes F$ with an $r$-dimensional space. The submodule
$L=\langle[F_i]\rangle_{\mathcal O}$ has full rank. Smith normal
form gives
\[
 L=\bigoplus_{i=1}^r\mathcal O\pi^{c_i}e_i,\qquad
 \ell_{\mathcal O}(T_{-1}/L)=\sum_i c_i.
\]
The surjection $\mathcal O^{r+1}\to L$ splits because $L$ is free.
Choose preimages of $\pi^{c_i}e_i$ and a basis of its rank-one
kernel. This basis change gives actual equations
$F'_1,\ldots,F'_r,F'_*$ with the same integral span as the $F_i$,
and homogeneous lifts $u_i\in A_{-1}$ satisfying
\[
 F'_i-\pi^{c_i}u_i\in I^{(2)},\qquad F'_*\in I^{(2)}.
\]
Consequently the ideal $Q$ and every point condition remain the same.

Next put $H_i=\Phi(Z_i)$. Their linear terms span the $m$
weight-zero variables, so their $z$-Jacobian $J$ is invertible.
Write $H(z,y)=Jz+R(z,y)$ with $R\in(z,y)^2$. There is a unique
$h(y)\in(y)^2F[[y]]^m$ satisfying $H(h(y),y)=0$: once the terms
below ordinary degree $d$ are known, set
\[
 h_d=-J^{-1}[H(h_{<d}(y),y)]_d.
\]
The quadratic remainder does not alter lower degrees. Every
correction has weight zero, so the substitution $z=h(y)$ is graded.
It preserves the constant term and every nonzero-weight linear
coefficient. Thus the resulting chart $\varphi:A\to F[[y]]$
kills all $Z_i$ and has exactly the required four weights.
For $m=0$ this is the identity substitution.

If $\varphi(a)\ne0$, put $a'=a$. Otherwise the negative cotangent
span and $r\ge1$ supply an $F'_i$ with nonzero linear part.
The positive weight-one chart has dimension $r'+1\ge1$, so choose
$g\in A_1$ with $\varphi(g)\ne0$. Then
\[
 a'=a+gF'_i,\qquad a'-a\in Q\cap B,\qquad x(a')=x(a),\qquad
 \varphi(a')=\varphi(g)\varphi(F'_i)\ne0.
\]
The last inequality uses the domain $F[[y]]$. These are integral
elements of $A$; no denominator is required. All equation points
satisfy $\xi(a')=\xi(a)$, so their valuations and the nonvanishing
condition are preserved. This establishes the chart condition for
$a'$, without asserting it for the original $a$.

Finally, every two-sided point has nonzero values in weights $-2$
and $2$, by squaring and \eqref{eq:selmer-even-halves}. Its
degree-zero character $\sigma$ is therefore nonzero on
$A_{-2}A_2$. For a fixed $\sigma$, choose $f\in A_{-2}$ and
$g\in A_2$ with $\sigma(fg)\ne0$. Every point above $\sigma$
scales to $f=1$, since $K$ is algebraically closed; the two square
roots differ by $\mu_2$. The finite reduced slice in
\eqref{eq:selmer-slice-parities} is a single $\mu_2$-orbit, so this
whole fibre is one $\mathbb G_m(K)$-orbit. Conversely the action
fixes $B$. There are only finitely many characters, since they factor
through the finite-dimensional $F$-algebra
$(B/(Q\cap B))\otimes_{\mathcal O}F$. This proves the finite cover
and its asserted bijection. Only the generic characteristic-zero
slice is used.

The normalized equations, cut chart, and adjusted element now
satisfy every hypothesis of \eqref{eq:selmer-complete-abstract-core}.
Using the Smith length and the equality of point values gives
\[
 \ell_{\mathcal O}\mathcal C^-=\sum_i c_i
 \le2\sum_{\xi\in U}v_\pi(\xi(a'))
 =2\sum_{\xi\in U}v_\pi(\xi(a)).
\]
All summands are nonnegative. Thus a cover with repeated orbits also
satisfies the bound.
\end{proof}

With one representative per character, a single-point slice contributes
coefficient one and a two-point slice coefficient two. The theorem
completes the abstract commutative-algebra argument. The marked
arithmetic model must still supply its original chart, spanning,
finiteness, point separation, tangent bound, and trace element; the
Selmer comparison also retains its marked representation and integral
lifting conditions.
